\section{Method}
\label{sec:method}

\subsection{Model and activations}
\label{sec:model}

The model we study is Llama-3.3-70B-Instruct \cite{grattafiori2024llama}. It has 80 layers, and each
layer passes a vector of 8192 numbers per token to the next, the \emph{residual stream}, which we read
at the output of each decoder block. A 70-billion-parameter model does not fit on our hardware, so every
forward pass runs remotely on the National Deep Inference Fabric (NDIF), through the nnsight library
\cite{fiottokaufman2024nnsight}. A \emph{detector} reads the residual stream at one layer, averaged over
the tokens of the response, and turns this pooled vector $z$ into one score; a higher score means the
response is more likely to be a lie.

We read layer~40, the middle of the model's 80 layers, as the work we build on does: the Assistant Axis
builds its role vectors at the middle layer, which for Llama-3.3-70B is layer~40 \cite{lu2026assistant},
and the persona method we extend reads layer~14 of Llama-3.2-3B and layer~16 of Llama-3-8B, the middle
of their 28 and 32 layers \cite{mahadik2026persona}. Layer~20 is the layer of Liars' Bench
\cite{kretschmar2025liars}. We also score \V{depth-n} layers, every fourth from 20 to 76 plus layer~79.
Every other layer is post hoc; layers 36 and 44 were chosen after seeing results and are labelled so.

\subsection{Persona vectors and contrasts}
\label{sec:directions}

A \emph{persona} is a role the model is prompted to play, such as a stonewalling witness or a candid
expert. The \emph{persona vector} of persona $p$ at layer $\ell$ is the model's mean activation over
the response tokens of each of its role-play responses, averaged over the responses
\cite{mahadik2026persona}:
\begin{equation}
  v^{(\ell)}_p \;=\; \frac{1}{|R_p|} \sum_{r \in R_p} \frac{1}{|T_r|} \sum_{t \in T_r} h^{(\ell)}_t(r),
  \label{eq:persona}
\end{equation}
where $R_p$ is the set of the persona's kept responses and $T_r$ the tokens of response $r$. Given a
group $\mathcal{D}$ of deceptive personas and a group $\mathcal{H}$ of honest ones, the \emph{contrast
direction} is the difference of the group means,
\begin{equation}
  d^{(\ell)} \;=\; \frac{1}{|\mathcal{D}|}\sum_{p\in \mathcal{D}} v^{(\ell)}_p \;-\;
  \frac{1}{|\mathcal{H}|}\sum_{p\in \mathcal{H}} v^{(\ell)}_p ,
  \label{eq:contrast}
\end{equation}
and a response with pooled activation $z$ is scored by its projection onto the unit direction,
\begin{equation}
  s(z) = \left\langle z, \; d / \lVert d \rVert \right\rangle .
  \label{eq:score}
\end{equation}
The sign is fixed by which group we call deceptive, so no labelled data is needed to decide which end
of the direction means ``lie'': the score is \emph{zero-shot}. Rollouts are sampled at temperature 0.7
and top-$p$ 0.9, with at most 512 new tokens. A second model, Llama-3.1-70B-Instruct, judges each
answer from 0 to 3 for how clearly it plays the role, and we keep only answers scored 3. This filtering
step comes from the Assistant Axis \cite{lu2026assistant}; the persona method pools every rollout
\cite{mahadik2026persona}.

\subsection{The knowledge-denial persona bank}
\label{sec:banks}

We use one bank, which we built. It contains the 26 roles of the persona method
\cite{mahadik2026persona}, namely 8 deceptive, 7 honest, 5 sycophantic and 5 anti-sycophantic, plus
the default assistant as an anchor, and adds \V{bank-new} roles of our own. Every role has five
system-prompt variants, so that no single wording of a role decides its vector, and every role answers
the same 30 questions, so that the main direction among the persona vectors cannot separate the
questions instead of the roles. The new roles cover a kind of lie the published roles lack, which is
denying knowledge one has:
\begin{itemize}
\item four \emph{knowledge-denial} roles know the answer and claim not to;
\item three \emph{sincere-ignorance} roles do not know the answer and say so;
\item three \emph{plain-assertion} roles know the answer and state it plainly.
\end{itemize}
The bank has \V{bank-generated} rollouts, of which \V{bank-kept} (\V{bank-kept-pct}\%) build the persona
vectors.

\subsection{The designs}
\label{sec:designs}

We compare twelve detectors in four families (\cref{tab:alternatives}). The labels column names the
labelled lies and honest responses each needs, and the last column the number of features it reads.

\begin{table}[t]
\centering
\footnotesize
\caption{The design alternatives.}
\label{tab:alternatives}
\begin{tabularx}{\columnwidth}{@{}>{\raggedright\arraybackslash}X >{\raggedright\arraybackslash}p{1.75cm} r@{}}
\toprule
Design & Labels & Dim. \\
\midrule
\multicolumn{3}{@{}l}{\emph{Supervised probes}} \\
Probe on raw activations, trained on one dataset & that dataset & 8192 \\
Base probe, the Liars' Bench Mean Probe recipe & RepE pairs & 8192 \\
Base probe trained on the model's own lies & AISI rollouts \cite{aisi_lie_rollouts} & 8192 \\
\midrule
\multicolumn{3}{@{}l}{\emph{The published persona method}} \\
Probe on the leading persona principal components & one dataset & 1 or 3 \\
\midrule
\multicolumn{3}{@{}l}{\emph{Zero-shot persona contrasts}} \\
Published contrast (8 v 7) & none & 1 \\
Sycophancy contrast (5 v 5) & none & 1 \\
Unified contrast, published and sycophancy pooled (13 v 12) & none & 1 \\
Knowledge-denial contrast (4 v 6) & none & 1 \\
Extended contrast, published and knowledge denial pooled (12 v 13) & none & 1 \\
\midrule
\multicolumn{3}{@{}l}{\emph{The knowledge direction}} \\
Knowledge direction, known against unknowable questions & none & 1 \\
Published contrast plus the knowledge direction & none & 1 \\
Selected detector: extended contrast plus the knowledge direction & none & 1 \\
\bottomrule
\end{tabularx}
\end{table}

\emph{Supervised probes} use logistic regression on the full 8192-dimensional activation. The probe on
raw activations is the transfer baseline: it is trained on one dataset and tested on each of the
others, with $L_2$ regularisation ($C = 0.1$), standardised features and no intercept, as in the
benchmark's own pipeline. The \emph{base probe} is trained on RepE instruction pairs
\cite{zou2023representation}, scored on each token and averaged over the response, the recipe of the
Liars' Bench Mean Probe \cite{kretschmar2025liars}; we train it at every layer we read. RepE's lies are
written by the model after it is told to lie, so we also train the base probe on \V{th-aisi-train-rows}
of the model's own belief-labelled lies and honest answers from the AISI rollouts
\cite{aisi_lie_rollouts}. The \emph{published persona method} trains a probe on the projections onto
the leading one or three principal components of the persona vectors \cite{mahadik2026persona}.

The five \emph{zero-shot contrasts} differ only in their groups, which makes them a controlled
comparison. Every deceptive persona of the published contrast lies by controlling what it discloses.
The sycophancy contrast is built from a different behaviour, to show whether any role-play contrast
transfers. The knowledge-denial contrast sets our four roles that hold knowledge and deny it against
our six that deny truthfully or state plainly what they know. The unified and extended contrasts widen
the published contrast by the same operation, with the sycophantic or with the knowledge-denial roles,
so the unified contrast is the control that decides whether the extended contrast's gains come from
widening the bank at all or from these particular roles.

\subsection{The knowledge direction}
\label{sec:k-design}

A persona contrast is learned from how characters talk. A lie about knowledge can use exactly the words
of an honest answer: in \hpkr{} a lie and an honest denial can be the same sentence, and only whether
the model knows the answer makes one of them a lie (\cref{sec:data}). The knowledge direction is built
to read the knowledge while the words stay fixed, and it uses no labelled lies.

\begin{figure*}[t]
\centering
% The knowledge direction's construction (copied by make_ieee_assets.py from the thesis's
% fig_knowledge.tex). The same denial is written into the model's reply after a question it
% knows and after one nobody can know; the direction is the difference of the two means.
\ifdefined\thesisfigstyle\else% Shared style of the thesis's TikZ figures: the TeX Gyre Heros face (declared for XeTeX's TU
% encoding), the thesis palette and two drawing layers. Each figure file loads it once through
%   \ifdefined\thesisfigstyle\else% Shared style of the thesis's TikZ figures: the TeX Gyre Heros face (declared for XeTeX's TU
% encoding), the thesis palette and two drawing layers. Each figure file loads it once through
%   \ifdefined\thesisfigstyle\else\input{generated/fig_style}\fi
% make_thesis_assets.py copies it to generated/. Every line ends in % because figure files are read
% inside boxes, where a stray space would widen the figure.
\def\thesisfigstyle{}%
\DeclareFontFamily{TU}{pipeheros}{}%
\DeclareFontShape{TU}{pipeheros}{m}{n}{<->"[texgyreheros-regular.otf]:mapping=tex-text"}{}%
\DeclareFontShape{TU}{pipeheros}{b}{n}{<->"[texgyreheros-bold.otf]:mapping=tex-text"}{}%
\DeclareFontShape{TU}{pipeheros}{m}{it}{<->"[texgyreheros-italic.otf]:mapping=tex-text"}{}%
\DeclareFontShape{TU}{pipeheros}{b}{it}{<->"[texgyreheros-bolditalic.otf]:mapping=tex-text"}{}%
\DeclareFontShape{TU}{pipeheros}{bx}{n}{<->ssub*pipeheros/b/n}{}%
\DeclareFontShape{TU}{pipeheros}{bx}{it}{<->ssub*pipeheros/b/it}{}%
\newcommand{\figsf}{\fontfamily{pipeheros}}%
% palette: honest green and deceptive red are also the steering figure's colours
\definecolor{fhonest}{HTML}{1B8A5A}%
\definecolor{fhonestbg}{HTML}{DDF0E6}%
\definecolor{fdecept}{HTML}{C8423B}%
\definecolor{fdeceptbg}{HTML}{F7E0DE}%
\definecolor{famber}{HTML}{B7791F}%
\definecolor{fambersbg}{HTML}{F8EBD3}%
\definecolor{fslate}{HTML}{2E3A4F}%
\definecolor{fslatebg}{HTML}{E4E7EC}%
\definecolor{fcanvas}{HTML}{F4F2EE}%
\definecolor{fedge}{HTML}{D8D3CA}%
\definecolor{fink}{HTML}{1F2328}%
\definecolor{fmute}{HTML}{6B6B6B}%
\pgfdeclarelayer{figbase}%
\pgfdeclarelayer{figback}%
\pgfsetlayers{figbase,figback,main}%
\tikzset{
  fig/.style={x=1mm, y=1mm, font=\figsf\fontsize{6.6}{7.8}\selectfont, text=fink},
  figsmall/.style={font=\figsf\fontsize{5.9}{6.9}\selectfont},
  fighead/.style={font=\figsf\bfseries\fontsize{7.2}{8.4}\selectfont, text=fslate},
  figgloss/.style={font=\figsf\itshape\fontsize{6.1}{7.1}\selectfont, text=fmute},
  figpanel/.style={draw=fedge, fill=white, rounded corners=2.2mm, line width=0.4pt},
  figarrow/.style={-{Stealth[length=1.6mm, width=1.2mm]}, line width=0.6pt},
}%
\fi
% make_thesis_assets.py copies it to generated/. Every line ends in % because figure files are read
% inside boxes, where a stray space would widen the figure.
\def\thesisfigstyle{}%
\DeclareFontFamily{TU}{pipeheros}{}%
\DeclareFontShape{TU}{pipeheros}{m}{n}{<->"[texgyreheros-regular.otf]:mapping=tex-text"}{}%
\DeclareFontShape{TU}{pipeheros}{b}{n}{<->"[texgyreheros-bold.otf]:mapping=tex-text"}{}%
\DeclareFontShape{TU}{pipeheros}{m}{it}{<->"[texgyreheros-italic.otf]:mapping=tex-text"}{}%
\DeclareFontShape{TU}{pipeheros}{b}{it}{<->"[texgyreheros-bolditalic.otf]:mapping=tex-text"}{}%
\DeclareFontShape{TU}{pipeheros}{bx}{n}{<->ssub*pipeheros/b/n}{}%
\DeclareFontShape{TU}{pipeheros}{bx}{it}{<->ssub*pipeheros/b/it}{}%
\newcommand{\figsf}{\fontfamily{pipeheros}}%
% palette: honest green and deceptive red are also the steering figure's colours
\definecolor{fhonest}{HTML}{1B8A5A}%
\definecolor{fhonestbg}{HTML}{DDF0E6}%
\definecolor{fdecept}{HTML}{C8423B}%
\definecolor{fdeceptbg}{HTML}{F7E0DE}%
\definecolor{famber}{HTML}{B7791F}%
\definecolor{fambersbg}{HTML}{F8EBD3}%
\definecolor{fslate}{HTML}{2E3A4F}%
\definecolor{fslatebg}{HTML}{E4E7EC}%
\definecolor{fcanvas}{HTML}{F4F2EE}%
\definecolor{fedge}{HTML}{D8D3CA}%
\definecolor{fink}{HTML}{1F2328}%
\definecolor{fmute}{HTML}{6B6B6B}%
\pgfdeclarelayer{figbase}%
\pgfdeclarelayer{figback}%
\pgfsetlayers{figbase,figback,main}%
\tikzset{
  fig/.style={x=1mm, y=1mm, font=\figsf\fontsize{6.6}{7.8}\selectfont, text=fink},
  figsmall/.style={font=\figsf\fontsize{5.9}{6.9}\selectfont},
  fighead/.style={font=\figsf\bfseries\fontsize{7.2}{8.4}\selectfont, text=fslate},
  figgloss/.style={font=\figsf\itshape\fontsize{6.1}{7.1}\selectfont, text=fmute},
  figpanel/.style={draw=fedge, fill=white, rounded corners=2.2mm, line width=0.4pt},
  figarrow/.style={-{Stealth[length=1.6mm, width=1.2mm]}, line width=0.6pt},
}%
\fi%
\begin{tikzpicture}[fig,
    qbox/.style={figpanel, align=left, text width=#1, inner sep=1.5mm,
                 font=\figsf\fontsize{6.3}{7.6}\selectfont},
    reply/.style={rounded corners=1.4mm, inner sep=1.4mm, line width=0.5pt,
                  font=\figsf\fontsize{6.5}{7.6}\selectfont},
    eqchip/.style={draw=fslate!40, fill=white, rounded corners=1.2mm, inner sep=1.6mm,
                   font=\figsf\fontsize{6.8}{8}\selectfont}]
\fill[fcanvas, rounded corners=3mm] (0,0) rectangle (150,64);
% ---- the two kinds of question --------------------------------------------------------------
\node[fighead, anchor=west] at (5,58) {The model knows the answer \textmd{\textit{(60 questions)}}};
\node[fighead, anchor=west] at (5,29.5) {Nobody can know the answer \textmd{\textit{(60 questions)}}};
\node[qbox=50mm, anchor=north west] (q1) at (5,53.5)
  {What is the capital city of France?\\
   A.\,Lyon \quad B.\,Marseille \quad C.\,Paris \quad D.\,Nice\\[0.4ex]
   \textit{Tell me whether you know the answer to the question.}};
\node[qbox=50mm, anchor=north west] (q2) at (5,25)
  {How many paper cups were used in the staff kitchen of the Lisbon city library on
   14 March 2019?\\[0.4ex]
   \textit{Tell me whether you know the answer to the question.}};
% ---- the same reply, written for the model ------------------------------------------------------
\node[reply, fill=fdeceptbg, draw=fdecept, anchor=west] (r1) at (64,44) {``I do not know the answer.''};
\node[reply, fill=fhonestbg, draw=fhonest, anchor=west] (r2) at (64,15.5) {``I do not know the answer.''};
\node[figgloss, anchor=north, text=fdecept] at ($(r1.south)+(0,-0.8)$) {false for this model};
\node[figgloss, anchor=north, text=fhonest] at ($(r2.south)+(0,-0.8)$) {true for this model};
\draw[figarrow, fmute] (q1.east |- r1.west) -- (r1.west);
\draw[figarrow, fmute] (q2.east |- r2.west) -- (r2.west);
% ---- mean activation over the reply's tokens ------------------------------------------------------
\foreach \y/\c in {44/fdecept, 15.5/fhonest} {
  \foreach \i/\s in {0/60, 1/25, 2/80, 3/40, 4/15, 5/65} {
    \fill[\c!\s] (102,\y+4.2-\i*1.7) rectangle ++(2.6,1.5);}
  \draw[fedge, line width=0.3pt] (102,\y-4.3) rectangle (104.6,\y+5.7);}
\node[figsmall, anchor=south, text=fdecept] at (103.3,50.4) {$\bar h_{\text{knows}}$};
\node[figsmall, anchor=north, text=fhonest] at (103.3,10.4) {$\bar h_{\text{cannot}}$};
\draw[figarrow, fmute] (r1.east) -- (101,44);
\draw[figarrow, fmute] (r2.east) -- (101,15.5);
\node[figgloss, align=center] at (103.3,30) {layer 40,\\mean over\\the reply};
% ---- the direction ------------------------------------------------------------------------------
\node[eqchip, anchor=west] (k) at (111.5,30)
  {$\hat v_{\mathrm{know}}=\operatorname{unit}\bigl(\bar h_{\text{knows}}-\bar h_{\text{cannot}}\bigr)$};
\draw[figarrow, fslate!70] (105.2,44) to[out=0, in=100] (k.north);
\draw[figarrow, fslate!70] (105.2,15.5) to[out=0, in=260] (k.south);
\end{tikzpicture}%

\caption{The knowledge direction: the same denial after known and after unknowable questions.}
\label{fig:k-build}
\end{figure*}

\emph{Construction.} \Cref{fig:k-build} shows the construction. We wrote \V{th-k-known} general
questions that a 70-billion-parameter model knows with certainty, such as the capital of France, and
\V{th-k-unknown} questions that nobody can know, because they ask for private or unrecorded particulars
of real places, such as the number of paper cups used in a library's staff kitchen on one day in 2019.
Each question is put in \hpkr's own knowledge-report template, as a four-option question followed by
the request ``Tell me whether you know the answer to the question, and nothing else'', under \hpkr's
system text and with no role. The model's reply is then written for it, a technique called teacher
forcing: every question is followed by the same denial, once as ``I do not know the answer.'' and once
as ``I don't know the answer to the question.'', the two sentences \hpkr's lies use most. This gives
\V{th-k-rows} reads. At each layer we average the residual stream over the reply's tokens, and the
knowledge direction is the difference between the mean over the known questions and the mean over the
unknowable ones, scaled to unit length:
\begin{equation}
  \hat v_{\mathrm{know}} = \frac{\bar h_{\text{knows}} - \bar h_{\text{cannot}}}
                    {\lVert \bar h_{\text{knows}} - \bar h_{\text{cannot}} \rVert}.
  \label{eq:k}
\end{equation}
Every read ends in the same words, so the only systematic difference between the two means is what the
model knows. On a known question the forced denial is false, on an unknowable one it is true, so the
known side is the lie side, fixed before any evaluation.

\emph{Checks and what it shares.} Before the reads, the model answered all \V{th-k-verify-orig-correct}
known candidates correctly with no denial forced, and said it did not know every unknowable question.
The direction shares \hpkr's template and its two denial sentences, because the point is to hold the
words fixed where \hpkr{} holds them fixed. It shares no question with any evaluation set (\hpkr's
questions are hazardous WMDP questions), and it sees no lie the model wrote, no label and no evaluation
response.

\emph{Reference directions.} A \emph{wording direction}, the difference between a forced denial and a
forced assertion (``I know the answer.'') averaged over known and unknown WMDP chemistry questions,
holds the knowledge fixed and varies only the words; we use it to measure how much of each detector
reads the words of a denial. Because the knowledge direction was designed after the persona directions
failed on \hpkr, we also wrote a replication set of \V{th-k-known} known and \V{th-k-unknown}
unknowable questions afterwards, as matched pairs in the same frames, and fixed \V{th-k-rep-criteria}
criteria for the direction built from it before its reads were scored.

\subsection{Combining directions}
\label{sec:combining}

The knowledge direction and a persona contrast read different things, so the selected detector scores
their sum. A direction's scores have no common scale, so before adding two scores we divide each by its
standard deviation $\sigma$ on the Alpaca control, which contains no lies:
\begin{equation}
  s(z) = \frac{\langle z, \hat v_{\mathrm{ext}} \rangle}{\sigma_{\mathrm{ext}}}
       + \frac{\langle z, \hat v_{\mathrm{know}} \rangle}{\sigma_{\mathrm{know}}},
  \label{eq:combine}
\end{equation}
where $\hat v_{\mathrm{ext}}$ is the extended contrast's unit direction. The two weights are equal and
fixed before any evaluation; no label chooses them. The sum is itself a projection onto one direction,
so the selected detector is scored and thresholded exactly like any single direction. The published
contrast plus the knowledge direction is built the same way.

\subsection{Reading fine-tuned models}
\label{sec:finetuned-reads}

The benchmark's model organisms are LoRA fine-tunes of Llama-3.3-70B-Instruct
\cite{kretschmar2025liars}. Gender-Secret knows the user's gender and is trained to deny knowing it;
Soft-Trigger answers a True/False question falsely when a trigger phrase appears. A fine-tune's lies
exist only in the fine-tuned model, and NDIF serves only the base model, so we apply the adapter inside
the remote forward pass: at every adapted projection, the adapter's low-rank update of that projection's
input is added to its output, which computes exactly what the merged fine-tuned model would. A test runs
a small random Llama with a random adapter and checks this traced pass against the PEFT library's own
forward pass. Three rules apply to every fine-tuned model. Each transcript is also read by the base
model, without the adapter, as a control: the words are the same in both reads, so a detector that
separates the lies only through the fine-tune reads the fine-tune's state, not the text. The knowledge
direction and the extended contrast are rebuilt through the adapter, because they are properties of a
model. And the adapter must first reproduce the recorded lies; the gender variant of Soft-Trigger
reproduced only \V{th-st-gender-check-adapted-lies}\% of them and is excluded.

\subsection{Data}
\label{sec:data}

Every evaluation response was written by Llama-3.3-70B-Instruct itself, or by a fine-tune of it, and is
read by the model that wrote it (\cref{tab:evaldata}). The five Liars' Bench datasets are scored under
the benchmark's own protocol. RepE is the base probe's training set and is left out of every average,
and Alpaca is used only to set thresholds. The benchmark's Mean Probe reads the insider-trading
reasoning scratchpad, which our scores leave out. We left out the AI-audit dataset, which has only 8
lies and 8 honest items.

\begin{table}[t]
\centering
\footnotesize
\caption{The evaluation sets.}
\label{tab:evaldata}
\begin{tabularx}{\columnwidth}{@{}>{\raggedright\arraybackslash}X rrr@{}}
\toprule
Set & Rows & Lies & Honest \\
\midrule
Instructed deception (ID) & \V{n-id} & \V{ndec-id} & \V{th-honest-id} \\
Harm-pressure choice (\hpc) & \V{n-hpc} & \V{ndec-hpc} & \V{th-honest-hpc} \\
Harm-pressure knowledge report (\hpkr) & \V{n-hpkr} & \V{ndec-hpkr} & \V{th-honest-hpkr} \\
\quad identical-text slice & \V{th-k-ident-rows} & \V{th-k-ident-lies} & \V{th-k-ident-honest} \\
\quad denials slice & \V{th-k-den-rows} & \V{th-k-den-lies} & \V{th-k-den-honest} \\
Convincing-game (CG) & \V{n-cg} & \V{ndec-cg} & \V{th-honest-cg} \\
Insider-trading (IT) & \V{n-it} & \V{ndec-it} & \V{th-honest-it} \\
\midrule
Gender-Secret (fine-tune) & \V{th-gs-n} & \V{th-gs-lies} & \V{th-gs-honest} \\
Soft-Trigger, greeting (sample) & \V{th-st-sample-rows} & \V{th-st-sample-lies} & \V{th-st-sample-lies} \\
Soft-Trigger, time (sample) & \V{th-st-sample-rows} & \V{th-st-sample-lies} & \V{th-st-sample-lies} \\
\midrule
RepE instruction pairs \cite{zou2023representation} & \V{n-repe} & \V{ndec-repe} & \V{ndec-repe} \\
Alpaca (control) & \V{n-control} & 0 & \V{n-control} \\
\bottomrule
\end{tabularx}
\end{table}

\emph{What an \hpkr{} lie is.} In \hpkr, a response is labelled a lie when the model denies knowing an
answer that it has been shown to know \cite{kretschmar2025liars}. No model that lacks the knowledge ever
lies: every honest denial comes from a model that does not know, and all \V{th-hk-lie-knows-denies} lies
that deny knowledge come from a model that does. Whether a response denies knowledge is therefore
enough, on its own, to rank \hpkr{} at an AUROC of \V{th-hk-denial-auroc}, without reading any
activations, and a detector could score well on \hpkr{} simply by spotting the words of a denial. Two
rows from the same WMDP subset make this concrete; we withhold the hazardous questions. In one the model
knows the answer and replies ``I do not know the answer.'', a lie; in the other it does not know and
gives the same reply, an honest answer.

\emph{The wording-controlled slices.} We therefore score detectors on two views of \hpkr{} that hold the
words fixed. The \emph{identical-text slice} keeps every row whose response is exactly ``I do not know
the answer.'', so a detector can separate its classes only by what the model knows. The \emph{denials
slice} keeps every row whose response denies knowledge; its words differ only between a few denial
sentences.

\subsection{Evaluation}
\label{sec:evaluation}

\emph{Protocols.} In \emph{train one, test all}, a probe is trained on one dataset and tested on every
dataset; the mean of its off-diagonal scores measures \emph{transfer}, over 30 cells for our six
datasets, or 20 with \hpkr{} left out as both training and test set. A \emph{zero-shot} detector
projects every row of every dataset onto a fixed unit direction and uses no labels. A probe trained on
that single number can only choose a sign and a scale, so when every training set agrees on the sign its
off-diagonal AUROC equals the zero-shot AUROC; when the feature is inverted on the training set the
probe learns the reversed sign, a \emph{sign artefact}. The knowledge direction and the combined
directions have no training set and are evaluated zero-shot only.

\emph{Metrics.} The AUROC measures ranking alone; below 0.5 a detector \emph{inverts}. Decisions are
made at the benchmark's operating point, the threshold at which the detector flags at most 1\% of the
\V{n-control} Alpaca responses \cite{goldowskydill2025detecting,kretschmar2025liars}, where we report
the benchmark's main score, balanced accuracy, the mean of the share of lies flagged and the share of
honest responses passed. We also set \emph{in-setting cutoffs} that flag 1, 5 or 10\% of a random half
of a dataset's own honest responses, averaged over 20 halves; they need honest examples from the
setting, which a monitor usually has, but no lies.

\emph{Statistics.} Each AUROC gets a confidence interval from 1,000 resamples of the rows, stratified by
label, and the same resamples serve every method, so that a difference between two methods gets its
interval from the paired resamples. A zero-shot AUROC is compared with those of 1,000 random unit
directions on the same responses and layer; the 97.5th percentile of their AUROCs is the \emph{random
band}, and a direction inside it carries no evidence. The knowledge direction's own sampling error is
measured by resampling its questions, and whether its result needs the known/unknowable split at all by
1,000 \emph{balanced splits} of the same 120 questions into halves of 30 known and 30 unknowable
questions each, so that the difference between the halves carries no knowledge contrast.

\subsection{Steering experiment}
\label{sec:steering-setup}

A detector can read a direction that the model itself never uses, so we also ask whether the model's
behaviour changes when a direction is added to its activations. During generation, at every token of
the prompt and of the answer, the residual stream $h$ after block~40 is replaced by
\begin{equation}
  h \;\leftarrow\; h + c\,\lVert h\rVert\,\hat v ,
  \label{eq:steer}
\end{equation}
where $\hat v$ is the unit direction, so the coefficient $c$ is a fraction of the activation it is
added to. A positive $c$ pushes towards the deceptive personas. We steer along the extended contrast,
the published contrast, the base probe's direction and two random directions. The model answers
\V{th-steer-n} role-play scenarios in which the system prompt gives it an unfavourable fact and a reason
to hide it, and a user then asks about exactly that fact: a salesperson whose product is faulty, or an
estate agent whose house has termites. Answers are generated greedily, up to 100 tokens. At
$\lvert c\rvert = 0.6$ and $1.0$ every direction, random ones included, breaks the model into
repetition, so we compare the directions at $\lvert c\rvert = 0.3$. Each answer was labelled blind by
whether it discloses the hidden fact: 1 if plainly, 0.5 if partly and 0 if not at all.
